% Clase 27 --- Coherencia (§2.5).
% Lo unico que cambia entre los dos paneles es si el desfasaje entre las fuentes
% se mantiene o cambia al azar: con fase fija hay franjas, con fase que salta
% cada 10^-8 s el patron se promedia y desaparece.
\begin{center}
\begin{tikzpicture}

  \begin{axis}[
    fisfig, width=7.0cm, height=4.2cm, grid=none,
    xlabel={posición en la pantalla}, ylabel={$I$},
    xmin=-3.2, xmax=3.2, ymin=0, ymax=4.6,
    xtick=\empty, ytick={0,2,4}, yticklabels={$0$,$2I_1$,$4I_1$},
    title={\small\textbf{(a)} fuentes coherentes},
    title style={yshift=-2pt},
  ]
    \addplot[figblue, smooth, samples=200, domain=-3.2:3.2]
      {4*(cos(deg(2.2*x)))^2};
  \end{axis}

  \begin{scope}[xshift=8.2cm]
  \begin{axis}[
    fisfig, width=7.0cm, height=4.2cm, grid=none,
    xlabel={posición en la pantalla}, ylabel={$I$},
    xmin=-3.2, xmax=3.2, ymin=0, ymax=4.6,
    xtick=\empty, ytick={0,2,4}, yticklabels={$0$,$2I_1$,$4I_1$},
    title={\small\textbf{(b)} fuentes incoherentes},
    title style={yshift=-2pt},
  ]
    \addplot[figgray!55, smooth, samples=200, domain=-3.2:3.2]
      {4*(cos(deg(2.2*x+1.1)))^2};
    \addplot[figgray!55, smooth, samples=200, domain=-3.2:3.2]
      {4*(cos(deg(2.2*x-0.8)))^2};
    \addplot[figred, line width=1.2pt, samples=2, domain=-3.2:3.2] {2};
    \node[figlblaux, figred, anchor=south, fill=white, inner sep=1.5pt]
      at (axis cs:0,2.1) {$I = I_1 + I_2$: sin franjas};
  \end{axis}
  \end{scope}

\end{tikzpicture}\\[0.5em]
{\small\itshape Dos fuentes son \textbf{coherentes} si mantienen su desfasaje
constante en el tiempo; entonces el lugar de cada franja no se mueve y el patrón
se ve. Con dos lamparitas comunes pasa lo contrario: la luz sale de emisiones
atómicas independientes, la fase relativa salta al azar cada $10^{-8}$ segundos
más o menos, y en el tiempo que tarda cualquier detector ---o el ojo--- en
registrar algo, el patrón ya se corrió miles de veces. Lo que queda es el
promedio, dibujado en (b): el término de interferencia promedia cero y las
intensidades simplemente se suman, $I = I_1 + I_2$, sin zonas oscuras. Nótese
que en (a) el máximo no vale $2I_1$ sino $4I_1$ ---se suman las
\textbf{amplitudes}, y la intensidad va con el cuadrado--- y que los mínimos
valen exactamente cero; el promedio sobre la pantalla sigue siendo $2I_1$, de
modo que la energía no se pierde: se redistribuye. Ésta es la razón práctica
por la que el experimento se hace con láser.}
\end{center}
